% 重绘自原书第 56 页的图 2.8；轴标签是原示例中的六个英文词元。
\begingroup
\colorlet{prefixcell}{blue!24}
\colorlet{crosscell}{green!32!gray!13}
\colorlet{targetcell}{yellow!37!orange!12}
\colorlet{maskedcell}{blue!12!gray!37}
\newcommand{\attentiontiles}[2]{%
  \node[font=\footnotesize\sffamily\bfseries] at (1.32,.49) {#2};
  \foreach \row/\ca/\cb/\cc/\cd/\ce/\cf in {#1} {
    \foreach \col/\tone in {0/\ca,1/\cb,2/\cc,3/\cd,4/\ce,5/\cf} {
      \node[tile,fill=\tone] at ({.53*\col},{-.53*\row}) {};
    }
  }
  \foreach \word [count=\idx] in {A,Survey,of,Large,Language,Models} {
    \node[anchor=east] at (-.24,{-.53*(\idx-1)}) {\word};
    \node[rotate=90,anchor=east] at ({.53*(\idx-1)},-2.91) {\word};
  }
}
\newcommand{\decoderbrackets}{%
  \draw[group] (-1.60,.17) -- (-1.73,.17) -- (-1.73,-2.82) -- (-1.60,-2.82);
  \node[rotate=90] at (-1.97,-1.32) {解码器};
  \draw[group] (-.19,-3.94) -- (-.19,-4.08) -- (2.83,-4.08) -- (2.83,-3.94);
  \node at (1.32,-4.29) {解码器};
}
\begin{tikzpicture}[
  font=\scriptsize\sffamily,
  tile/.style={draw=black!56,minimum size=4.1mm,inner sep=0pt,rounded corners=1pt,line width=.48pt},
  group/.style={draw=black!49,line width=.52pt,rounded corners=1pt}]
\begin{scope}[shift={(2.74,5.13)}]
\attentiontiles{
  0/prefixcell/maskedcell/maskedcell/maskedcell/maskedcell/maskedcell,
  1/prefixcell/prefixcell/maskedcell/maskedcell/maskedcell/maskedcell,
  2/prefixcell/prefixcell/prefixcell/maskedcell/maskedcell/maskedcell,
  3/crosscell/crosscell/crosscell/targetcell/maskedcell/maskedcell,
  4/crosscell/crosscell/crosscell/targetcell/targetcell/maskedcell,
  5/crosscell/crosscell/crosscell/targetcell/targetcell/targetcell%
}{因果解码器}
\decoderbrackets
\end{scope}
\begin{scope}
\attentiontiles{
  0/prefixcell/prefixcell/prefixcell/maskedcell/maskedcell/maskedcell,
  1/prefixcell/prefixcell/prefixcell/maskedcell/maskedcell/maskedcell,
  2/prefixcell/prefixcell/prefixcell/maskedcell/maskedcell/maskedcell,
  3/crosscell/crosscell/crosscell/targetcell/maskedcell/maskedcell,
  4/crosscell/crosscell/crosscell/targetcell/targetcell/maskedcell,
  5/crosscell/crosscell/crosscell/targetcell/targetcell/targetcell%
}{前缀解码器}
\decoderbrackets
\end{scope}
\begin{scope}[xshift=5.50cm]
\attentiontiles{
  0/prefixcell/prefixcell/prefixcell/maskedcell/maskedcell/maskedcell,
  1/prefixcell/prefixcell/prefixcell/maskedcell/maskedcell/maskedcell,
  2/prefixcell/prefixcell/prefixcell/maskedcell/maskedcell/maskedcell,
  3/crosscell/crosscell/crosscell/targetcell/maskedcell/maskedcell,
  4/crosscell/crosscell/crosscell/targetcell/targetcell/maskedcell,
  5/crosscell/crosscell/crosscell/targetcell/targetcell/targetcell%
}{编码器—解码器}
\draw[group] (-1.60,.17) -- (-1.73,.17) -- (-1.73,-1.26) -- (-1.60,-1.26);
\draw[group] (-1.60,-1.38) -- (-1.73,-1.38) -- (-1.73,-2.82) -- (-1.60,-2.82);
\node[rotate=90] at (-1.97,-.55) {编码器};
\node[rotate=90] at (-1.97,-2.10) {解码器};
\draw[group] (-.19,-3.94) -- (-.19,-4.08) -- (1.26,-4.08) -- (1.26,-3.94);
\draw[group] (1.38,-3.94) -- (1.38,-4.08) -- (2.83,-4.08) -- (2.83,-3.94);
\node at (.53,-4.29) {编码器};
\node at (2.12,-4.29) {解码器};
\end{scope}
\end{tikzpicture}
\endgroup
